\section{Build \& Self-Verify：让 Agent 学会自我检验}

\subsection{问题：模型的"第一印象偏见"}

当前模型存在一个显著的行为偏见：\textbf{倾向于接受自己写出的第一个看似合理的解决方案}。Trace 分析中最常见的失败模式就是——Agent 写了一个方案，重新读了一遍自己的代码，认为"看起来没问题"，然后直接提交。

这种行为缺少了软件工程中至关重要的一步：\textbf{测试}。

\begin{importantbox}{Self-Verification 的核心洞察}
今天的模型是出色的"自我改进机器"（self-improvement machines）。Self-Verification 允许 Agent 通过运行时反馈在单次执行中自我改进。但模型天然\textbf{不会}主动进入 Build-Verify 循环——需要 Harness 显式引导。
\end{importantbox}

\subsection{四阶段问题解决框架}

LangChain 在 System Prompt 中注入了一个结构化的问题解决流程：

\begin{enumerate}[leftmargin=2em]
    \item \textbf{Planning \& Discovery（规划与探索）}：阅读任务描述，扫描代码库，基于任务规格和验证方式制定初始计划
    \item \textbf{Build（构建）}：按计划实现，同时构建测试——不仅测试 Happy Path，还要覆盖 Edge Cases
    \item \textbf{Verify（验证）}：运行测试，完整阅读输出，\textbf{对照任务规格}（而非对照自己的代码）进行验证
    \item \textbf{Fix（修复）}：分析错误，回溯原始规格，修复问题
\end{enumerate}

\begin{warningbox}{Verify 阶段的关键区分}
验证时必须\textbf{对照任务规格（Task Spec）}，而非对照自己写的代码。Agent 常犯的错误是：对照自己的代码验证自己的输出——这是循环论证，无法发现偏离需求的问题。
\end{warningbox}

\subsection{PreCompletionChecklistMiddleware}

除了 Prompt 引导，LangChain 还引入了一个确定性机制：\textbf{PreCompletionChecklistMiddleware}。这个中间件在 Agent 准备退出时拦截它，强制注入一个提醒，要求 Agent 对照 Task Spec 做一轮完整的验证。

这类似于所谓的 Ralph Wiggum Loop——通过 Hook 机制在 Agent 即将退出时强制其继续执行验证步骤。

\begin{knowledgebox}{确定性 vs 概率性控制}
Harness 设计中存在两种控制策略：
\begin{itemize}[leftmargin=1.5em]
    \item \textbf{概率性控制}：通过 Prompt 引导模型行为（模型可能遵循也可能忽略）
    \item \textbf{确定性控制}：通过 Middleware/Hook 强制执行（代码层面保证一定触发）
\end{itemize}
两者互补：Prompt 设定意图方向，Middleware 兜底确保关键步骤不被跳过。LangChain 在 Self-Verification 上同时使用了这两种策略。
\end{knowledgebox}

\subsection{本章小结}

Self-Verification 是 LangChain 在 Harness 改进中收益最大的单项优化。核心是通过 Prompt 引导和 Middleware 兜底，迫使 Agent 从"写完就交"变为"写完-测试-对照规格验证-修复"的闭环。关键在于验证对象是\textbf{任务规格}而非 Agent 自身代码。


